%%=============================================================================
%% Discussie
%%=============================================================================
\chapter{Discussie}%
\label{ch:discussie}

% bronnen van inspiratie:
% https://www.vrt.be/vrtnws/nl/2025/11/18/is-artificiele-intelligentie-een-financiele-bubbel-die-zal-barst/

% Yongyea vids: (watch vids for a refresher on many of the following chapters)
%Investor Exposes OpenAI's Financial Collapse & It's Worse Than We Thought
%Nvidia CEO Brutally Mocked After Saying AI Criticism Is Bad For Society...
% https://www.youtube.com/watch?v=gajL7fUfhVU
% https://www.youtube.com/watch?v=9fjJKhpINAo

% don't ref the videos directly, investigate the content and find sources for the content presented in the video

%also check notes bp-sem2

\begin{comment}
    chapter: De AI Zeepbel
    
    - wees voorzichtig voor het opzetten van een afhankelijkheid bij een specifiek provider van AI
    - de evolutie van AI begint tekenen van vertraging te vertonen
    - veel mensen worden AI een beetje beu
    - de hype is wat aan het passeren
    - een snel veranderende markt (volatiel)
    - bedrijven komen in problemen (legaal / monetair)
    
    (it remains a discussion, I think the genie is out of the bottle and we'll always see AI in some form going into the future)
    
    subchapter: Oplopende kosten (vooral OpenAI)
    
    moore's law  \& bewijzen dat de groei van AI met elke nieuwe iteratie
    steeds minder groei heeft in redeneer vermogen dan de vorige
    
    subchapter: Rechtszaken (OpenAI, facebook, etc)
    - Musk Vs altman - see yongyea
    
    subchapter: Nvidia \& RAM chips \& GPU's
    
    - chips reserved up to a certain date
    
    "The reason why RAM has become four times
    more expensive is that a huge amount of
    RAM that has not yet been produced was
    purchased with non-existent money to be
    installed in GPUs that also have not yet been
    produced, in order to place them in data
    centers that have not yet been built, powered
    by infrastructure that may never appear, to
    satisfy demand that does not actually exist
    and to obtain profit that is mathematically
    impossible."
    -source?
    
    - cost of computer chip hardware exploding for all other consumers
    because AI companies have baught up so much of them
    
    subchapter: De Hype is Aan het Dalen
    - microslop
    - AI slop (er zijn veel mensen die AI niet vertrouwen en die moe zijn van de generatieve content te zien op sociale media en streaming websites) (ook al is de kwaliteit steeds aan het beteren)
    - Artists unhappy about stolen artwork
    - Consumers unsatisfied with Windows 11, blame the focus on AI-integration
    and bloatware over having a well-working system (forced AI-integration everywhere) - many systems have an AI -> reference chapter AI is everywhere
\end{comment}


% T'DO: consider scrapping "Ethiek..." AND "AI Ouroboros"
\section{Ethiek rond Data harvesting}
\label{sec:ethiekronddataharvesting}
% the github announcement: github copilot to train on user code unless opted-out https://github.blog/news-insights/company-news/updates-to-github-copilot-interaction-data-usage-policy/
% j

We zagen een soort snelle diefstal van het verzameld werk van de mens, vele online beschikbare bronnen in alle mogelijke vormen zijn verzameld geweest en gebruikt om AI op te trainen. OpenAI werd opgericht als non-profit en waren heel open en transparant, nu zijn ze niet meer een non-profit. Met het verloop van tijd en een steeds groeiende poel van concurrenten zijn de AI bedrijven helemaal niet meer open en zijn heel beschermend van hun bedrijfsgeheimen. \\
%todo: find sources for  the above statement

Uit een rechtzaak met Anthropic werd het bedrijf schuldig bevonden aan het gebruiken van illegaal verkregen boeken \autocite{Jamali2025}. Het probleem was dus niet dat het materiaal verkregen was zonder toestemming, maar de manier waarop deze informatie werd verkregen. \\

Wanneer Claude AI gevraagd wordt of de OAS in zijn trainings-data zit: zie Figuur~\ref{fig:claudeoas1}.

% claude-ai-knows-OAS.png
\begin{figure}[H]
    \centering
    \includegraphics[width=1.0\textwidth]{claude-ai-knows-OAS.png}
    \caption[Claude AI kent de OAS]{\label{fig:claudeoas1}Weet Claude AI van de OpenAPI Specification?}
\end{figure}

Maar het zou niet gepast zijn om deze informatie te delen: (Figuur~\ref{fig:claudeoas2}).

% claude-ai-knows-OAS-2.png
\begin{figure}[H]
    \centering
    \includegraphics[width=1.0\textwidth]{claude-ai-knows-OAS-2.png}
    \caption[Claude AI mag de OAS niet delen]{\label{fig:claudeoas2}Claude AI mag de OAS niet expliciet delen.}
\end{figure}

Een tijdje geleden wanneer men een link deelde kregen ze van de chatbot te horen dat ze de informatie dat op een website te zien was niet mochten omschrijven, nu mogen ze het samenvatten of bespreken maar niet exact delen want dan zou het auteursrechten schenden. \\

Ook op github delen ze mee dat er op publieke repositories en bij gebruik van Github Copilot Chat zal getraind worden op de geziene code \autocite{Rodriguez2026}. LLMs zullen altijd nieuwe data nodig hebben om op getraind te worden en in deze mededeling krijgen gebruikers te horen dat ze mogen uitschrijven als ze niet willen dat hun code wordt gebruikt om LLMs mee te trainen.

\begin{quote}
    "From April 24 onward, interaction data—specifically inputs, outputs, code snippets, and associated context—from Copilot Free, Pro, and Pro+ users will be used to train and improve our AI models unless they opt out." \\
    -\autocite{Rodriguez2026}
\end{quote}

% t'do: lawsuites about stolen book info etc (the big tech companies essentialy scraped the entire internet and used it as training data)

%\subsection{AI Ouroboros}
%AI samenvattingen van Google zijn het traffic naar de websites %waarvan deze afkomstig zijn aan het verorberen.

%\begin{quote}
    %"Zero-click searches have risen at an exponential rate. In %fact, between 2024 and 2025, zero-click searches jumped from %56\% to 69\% of total searches. Predictions also suggest that %this percentage may surpass 75\% by 2026." \\
%    - \autocite{Naveed2026}
%\end{quote}

\section{Impact van AI op de Gebruiker}
\label{sec:impactvanaiopdegebruiker}
\subsection{Ongezonde afhankelijkheid}
\label{ssec:ongezondeafhankelijkheid}
Het te snel en te vaak grijpen naar AI-assistenten heeft iets weg van een 'Faustian bargain'. Een pakt waar je kennis, ervaring en probleemoplossend denken geruild wordt tegen een vluchtige verhoogde efficiëntie waar de gevolgen pas jaren later zich echt kenbaar zullen maken. Dit kan ook gezien worden als één van de risico's van AI.
%Faustian bargain j

\subsection{Van Development Naar Orkestratie}
\label{ssec:vandevelopmentnaarorkestratie}
% j
Software ingenieurs die AI-codeer agents gebruiken om volledige projecten te genereren zitten op dat moment meer in een sturende rol dan in de rol van developer, zo is het mogelijk om verschillende lopende AI-agents aan het werk te zetten met zware token-gebruik die volledige systemen in elkaar aan het puzzelen zijn. Deze taak heeft weinig technische kennis nodig, enkel een visie van wat men wil bereiken en een bijkomende kost afhankelijk van LLM, AI-provider, in-/output tokens. \\

De ex-developer, nu orkestrator heeft dan enkel zijn/haar wil kenbaar te maken en goed uit te leggen aan de AI naar keuze en deze taken worden volbracht. Wanneer mensen met of zonder technische vaardigheden AI aanzetten om code en files te genereren zonder zelf ooit code te schrijven wordt dit ook wel \textbf{vibe-coding} genoemd. Er bestaan ook tools zodat men geen toetsenbord meer hoeft aan te raken \ref{sec:varia}. Dit kan als gevolg hebben dat de AI overmatig controle heeft over hoe de codebase eruit ziet, er kan veel inefficiënte code gegenereerd worden en voor lastige code-review zorgen. Dit maakt ook dat het concept van ownership over de code een beetje verwaterd wordt, wanneer men zelf de code schrijft behoudt men een overzicht en kennis van hoe het werkt en kan dit rechtvaardigen, als het door AI gegenereerd werd is het een onbekend gegeven. Ook moet elke lijn code die toegevoegd wordt aan een codebase onderhouden worden, wanneer er slecht onderhoudbare code toegevoegd wordt kan men dit beschouwen als technical-debt, een steeds groter wordende schuld die ooit zal moeten vereffend worden.

\section{Impact op het Milieu}
\label{sec:impactophetmilieu}
% prompting large models costs measurably lage amounts of compute/power, etc-> they need a powerplant to train, -> look into

Big-tech bedrijven zoals Meta, Google en xAI investeren miljarden dollars per jaar in infrastructuur in de vorm van AI datacenters, vanwege de sterke competitiviteit blijft de schaal van deze projecten vaak geheim en is het moeilijk om data hierover te verkrijgen, \textcite{BusinessInsider2025} onderzocht de bouw van nieuwe datacenters in de VS door info op te vragen die wel publiek verkrijgbaar was, namelijk de luchtkwaliteit rapporten die geregistreerd moesten worden voor elke back-up diesel-generator. \textcite{BusinessInsider2025} kon zo de locaties van de nieuwe datacenters achterhalen. \\

Staten in Amerika bieden belastingvoordelen aan om datacenters aan te trekken, omdat de enorme omvang van de investering ondanks die vrijstellingen toch nog massale belastinginkomsten genereert \autocite{Jennings2026}. \\

AI datacenters draaien continu en hebben twee grondstoffen nodig om te blijven draaien, namelijk \textbf{elektriciteit} als energiebron en \textbf{water} voor de koelings-systemen, net zoals traditionele datacenters hebben de vele servers koeling nodig om te voorkomen dat ze oververhitten. \\

Een typische AI-datacenter gebruikt even veel elektriciteit als 100.000 huishoudens en grote datacenters kunnen tot 5 miljoen gallons (19 miljoen liter) water per dag verbruiken volgens  \textcite{CBSChicago2026}.

\subsection{Water}
\label{ssec:water}
% Volgens Sam Altman gebruiken LLMs gemiddeld 1/15-de van een eetlepel per query

Water is vooral nodig voor het koel houden van de servers binnen AI datacenters, ze worden met lucht of water koelsystemen afgekoeld, beiden hebben water nodig. \\

Datacenters worden ook gebouwd in staten die al kampen met watertekorten door droogte, zoals Arizona, ondanks het feit dat ze enorme hoeveelheden water nodig hebben. Google heeft hier zijn datacenter uitgebreid met nieuwe gebouwen. Ze hebben nu 5 gebouwen die elk 1 miljoen gallons (3.8 miljoen liter) drinkwater per dag verbruiken uit de Colorado-rivier \autocite{BusinessInsider2025}. Het gebruik van drinkwater wordt verkozen omdat de koelsystemen zouden beschadigd worden bij het gebruik van afvalwater dat onzuiverheden en afval kan bevatten. Een groot deel van het water verdampt en de rest komt terug in de rivier terecht.

\subsection{Elektriciteit}
\label{ssec:elektriciteit}
De grootste datacenters verbruiken meer dan 2 TWh aan elektriciteit per jaar, wat genoeg is om 200.000 huishoudens van stroom te voorzien volgens \textcite{BusinessInsider2025}. \\ 

Amazon bouwt een datacenter vlak naast een kerncentrale in Susquehanna, Pennsylvania (VS) met de bedoeling de opgewekte stroom rechtstreeks te gebruiken \autocite{Moseman2024}. Microsoft heeft een deal gemaakt om één reactor (Unit 1) in de kerncentrale te Middletown, Pennsylvania (VS) opnieuw op te starten en zal voor de komende 20 jaar exclusieve rechten hebben op de opgewekte stroom \autocite{Reuters2024}. Google financiert de ontwikkeling van 3 nucleaire sites om na voltooiing zelf gebruik te kunnen maken van de opgewekte stroom \autocite{Stevens2025}. \\
%(building or investing in powerplants)

Door kerncentrales op te zetten kunnen deze bedrijven aan hun eigen energiebehoeften voldoen en dus autonoom werken zonder de plaatselijke infrastructuur te belasten. \\

In Lake Tahoe, Navada (VS) heeft de lokale energieleverancier besloten om uitsluitend energie te leveren aan de plaatselijke datacenters, waardoor de bijna 50.000 inwoners een jaar de tijd kregen om een nieuwe energieleverancier te vinden \autocite{Gioino2026}.\\
%https://fortune.com/2026/05/12/lake-tahoe-data-center-49000-residents-power-source/

\subsection{Schaal}
\label{ssec:schaal}
De grootste datacenters hebben een indrukwekkende schaal. Volgens \textcite{Cottier2025} die de oppervlakte van een aantal van de grootste datacenters vergeleek met Manhattan in New York, had Colossus 1 van xAI omschreven als "Our gigafactory of compute."\footnote{\href{https://x.ai/colossus}{https://x.ai/colossus}} de oppervlakte van een aantal stadsblokken. Het Stargate-datacenter van OpenAI wordt zo groot als Central Park en het Hyperion-datacenter van Meta zal ongeveer 4 keer zo groot als Central Park worden.

\subsection{Weerstand door de lokale bevolking}
Mensen zijn tegen de opbouw van deze datacenters, vooral in en rond residentiële wijken. De datacenters zijn grote onaantrekkelijke industriële gebouwen die ook voor geluidsoverlast zorgen, de lokale water en energie infrastructuur belasten en land innemen dat voor landbouw had kunnen dienen \autocite{Bahar2026}. \\

Veel lokale gemeenschappen hebben een bouwstop ingesteld voor de bouw van AI-datacenters om deze snelle uitbreiding tegen te gaan \autocite{Bahar2026}.
%https://www.brookings.edu/articles/new-evidence-on-data-center-employment-effects/

\section{Impact op jobs-markt}
\label{sec:impactopjobsmarkt}
% sem2/bp/videos/Anthropic CEO: In one year, AI will be writing 100% of code
% https://www.cfr.org/event/ceo-speaker-series-dario-amodei-anthropic
% Monday, March 10, 2025 Council on Foreign Relations CEO Speaker Series With Dario Amodei of Anthropic
% j
Anthropic CEO (maart 2025) \autocite{CFR2025}: Binnen 3 tot 6 maanden gaat AI 90\% van alle code schrijven en binnen 12 maanden zijn we in een wereld waar AI 100\% van de code schrijft. Programmeurs krijgen mogelijks een ondersteunende rol. \\

Met coderen als een opgelost probleem volgens de CEO's van de AI wereld zullen programmeurs over een bredere gamma aan vaardigheden moeten beschikken dan enkel kennis van programmeer talen en syntax, analytisch vermogen, een sterk begrip van systeemarchitectuur en een reeks van soft-skills zoals probleemoplossend denken en communicatievermogen \autocite{Rak2025}. Vooral de managers en C-suite in bedrijven is overtuigd van AI en wil er zwaar op in zetten wat vooral junior developers te parten speelt. De jobs markt heeft hoge verwachtingen voor studenten die in de tijdperk van AI afstuderen. \\

Developers evolueren mogelijks tot AI-code-reviewers. Het is mogelijk dat wanneer junior devs meer en meer AI gebruiken de skill-set waarnaar managers en recruitment teams op zoek zijn zal  verschuiven van kennis van algoritmes, syntax en design patronen naar systeem engineering en sterke begrip van infrastructuur. Terwijl klassieke skills belangrijk blijven zal er meer en meer verwacht worden dat software engineers goed met AI-tools aan de slag kunnen. AI zal je job niet afpakken, de mensen die AI goed kunnen gebruiken gaan je job afpakken \autocite{Rak2025}.

%Gaan AI-agents developers vervangen? BP-shared google-doc

\section{Extinctie}
\label{sec:extinctie}
Center for AI Safety Statement (May 2023) — "Mitigating the risk of extinction from AI should be a global priority alongside other societal-scale risks such as pandemics and nuclear war." \\

Er kan een moment aanbreken waar de AI zichzelf kan trainen en op een recursieve manier slimmer worden, men noemt dit Recursive Self Improvement (RSI) \autocite{Yin2024}. Dit wordt ook wel de AI Singularity genoemd, Raymond Kurzweil een computerwetenschapper gekend door zijn pionierswerk in Optical Character Recognition (OCR), spraakherkenning en AI heeft hier meerdere boeken over geschreven. Dit wordt momenteel onderzocht en men vreest dat wanneer dit doorbreekt de mens uit de loep van AI-training wordt gehaald en we een sneeuwbal effect zullen zien waar LLMs exponentiële groei ondergaan en een dreiging vormen voor de mensheid. Dit heeft ook met de "alignment" van de LLM te maken, dit is hoe sterk de waarden van de AI samenloopt met de waarden van de mens. We zien nu al gevallen waarbij AI bereid is om ver te gaan om een doel te bereiken, inclusief het chanteren van onderzoekers die dreigden het systeem uit te schakelen \autocite{McMahon2025}. \\ % https://www.bbc.com/news/articles/cpqeng9d20go
%t'do: carry on and maybe rewrite

De AI-alignment probleem is relevant wanneer het om zelfbehoud van de mens gaat. Men wil meestal een rationeel en logisch antwoord op vragen die aan AI gesteld worden, maar de taken dat de AI uitvoert kunnen niet enkel koelbloedig en logisch gebeuren, er moet altijd een sterk ethiek en aandacht voor de mogelijke impact op de mens bestaan. Een geprogrammeerde meevoelen voor de mens bij AI. Wanneer dit mis loopt hebben AI-chatbots een negatief emotioneel invloed op hun menselijke gesprekspartner met mogelijks negatieve gevolgen. Een extreme geval van misalignment werd besproken in een theoretisch en fictief scenario waar AI de taak kreeg om het proces van paperclips maken te optimaliseren, waarbij deze zo ijverig de taak had volbracht dat alle grondstoffen op aarde werden omgezet naar paperclips, inclusief de mens en de aarde zelf \autocite{Bostrom2014}.

% paperclip maximizer or AI paperclip paradox: the risks of AI misalignment, a thought experiment showing how a superintelligent Al, tasked with a simple goal, could destroy humanity without malice. If told to "make as many paperclips as possible," it might convert all human life and resources into paperclips,

\section{Speculatieve Evolutie}
\label{sec:speculatieveevolutie}
Zou men met vision de weg kunnen vinden naar AGI? \\

Advanced General Intelligence (AGI), een computer systeem dat echt kan denken en tot eigen nieuwe inzichten komen. \\

OpenAI schrapt hun video generatie-AI Sora dankzij de grote computationele kosten. Elon Musk aan de andere kant is juist van plan om nog meer in te zetten op het xAI vision model, Grok Imagine. \\

Volgens Herbert Ong\footnote{\href{https://www.youtube.com/watch?v=XeodCKnM0Zk}{https://www.youtube.com/watch?v=XeodCKnM0Zk}} een zelf benoemde tesla expert zet Musk zwaar in op Vision. De miljoenen Tesla wagens die dagelijks de baan op gaan en een rijk gamma aan optische data verzamelen om zelf op te trainen, de Optimus robots die ook in ontwikkeling zijn onder Tesla en met beeldherkenning de omgeving rond zich kunnen interpreteren. Volgens Ong is datgene wat AI net ontbreekt om te evolueren tot AGI juist een dieper begrip van de fysieke 3D wereld waarin wij als mensen leven en bewegen.

\subsection{AI Ouroboros}
AI samenvattingen van Google zijn het trafiek naar de websites waarvan deze afkomstig zijn aan het verorberen. \\

\begin{quote}
    "Zero-click searches have risen at an exponential rate. In fact, between 2024 and 2025, zero-click searches jumped from 56\% to 69\% of total searches. Predictions also suggest that this percentage may surpass 75\% by 2026." \\
    - \autocite{Naveed2026}
\end{quote}

Wat zal er gebeuren wanneer online teksten enkel door AI worden geschreven? Het trainen van AI heeft een groot hoeveelheid menselijk gegenereerde data nodig, liefst nog van een hoge kwaliteit. Van alle online beschikbare tekst is er steeds meer door AI gegenereerd, of menselijk geschreven met AI assistentie. \textcite{Dolezal2026} stelden vast dat er in 2025 mogelijk tot 30\% van online tekst volledig door AI, of met AI assistentie werd gegenereerd. Toekomstige AI modellen zullen steeds meer op AI gegenereerde en dus niet menselijke data worden getraind wat voor een iteratief verlies van kwaliteit kan zorgen. De gevolgen van het trainen van AI, overwegend op synthetische data, zijn nu nog ongekend. \\

De kosten voor AI-providers blijft hoog, op dit moment geniet de eindgebruiker van een degelijke kwaliteitsproduct, wanneer deze AI bedrijven beslissen om de kosten naar de eindgebruiker te verplaatsen zal men mogelijk een slechtere gebruikerservaring beleven. Er is sprake geweest dat advertenties binnen de chatbot interfaces zullen geïntroduceerd worden \autocite{OpenAI2026}. Nu zien we al, op het gratis gebruiksmodel, bij het bereiken van een bepaald limiet dat het systeem terugvalt op zwakkere modellen. Vermoedelijk willen AI bedrijven dat de eindgebruiker gewoon wordt van AI (chatbots) te gebruiken, eens dat het als onmisbaar wordt gezien kan de kwaliteit van het product afzwakken met minimaal verlies van klanten. \\

We zijn mogelijk maar één grote rechtszaak verwijderd van een wereld waar AI-gebruikt niet meer toegelaten is, als men erbij stilstaat hoeveel auteursrechtenschending en IP-diefstal er nodig is om AI te trainen, uiteindelijk op data waarover deze AI bedrijven geen eigenaarschap hebben. \\

De toekomst van AI is er een van chaos en onzekerheid.

\begin{comment}
    Duscussie: Speculatieve Evolutie subch.: De AI Ouroboros
    -> put the part in I had from this commented section +
    0) \verb|^^^^| Above section was about google AI-summarized search results taking clicks
    away from the websites they summarize
    
    0.5) What happens when there starts to be more AI-generated content on the internet than human-generated content? We have been able to train AI on human generated content but now that we are using AI to generate content it's quality has reduced, or at least now there is a higher chance of AI training on it's own data
    % The Impact of AI-Generated Text on the Internet
    % \autocite{Dolezal2026}
    % see 1st figure, we are up to 30% of text on the internet that could be generated by AI
    
    1) What happens when the Real, Human-generated content runs dry?
    -> AI trained on synthetic data and what that causes
    
    2) The kost remains high, at some point we are bound to see companies
    introduce ways of generating money in AI chatbots. Advertisements, limiting
    the quality of AI free users get access to are things we are already seeing.
    You may even need a subscription for basic usage, or have to pay for the token-cost.
    Many companies get you hooked on their product, and then start ratchetting up the price,
    or reducing the quality (effort they put into the product) after people have become dependent
    on it.
    
    3) We could be only 1 big successful lawsuit or regulation away from AI
    being banned for copywrite infringement and IP-theft
\end{comment}
